\chapter*{Lời mở đầu}
\addcontentsline{toc}{chapter}{Lời mở đầu}
\label{ch:loi-mo-dau}

\section*{1. Lý do chọn đề tài}
Caro 11×11 có luật đủ đơn giản để kiểm chứng bằng một engine dùng chung, đồng thời có đủ tình huống chiến thuật để so sánh cách ra quyết định. Đồ án xây dựng tám cấp tác tử, từ chọn nước ngẫu nhiên, luật, tìm kiếm đến policy học được và tự đấu với PUCT. Các đối sánh đổi màu trên cùng khai cuộc giúp phân biệt sức chơi với một phần ảnh hưởng của bên đi trước.

Các kết quả chính cho thấy cắt tỉa AlphaBeta giảm chi phí trên bộ thế đã benchmark.
Với cùng mạng L5, thêm PUCT ở 200 mô phỏng nâng điểm trước AlphaBeta từ 21,0\% lên
58,5\%. Checkpoint L8 được giữ ở vòng 37 đạt 82,5\% trước L7 và 68,75\% trước AlphaBeta
trong hai lượt đánh giá riêng. Kết quả đủ để so sánh trong các thiết lập đã đo,
còn sức chơi theo cùng thời gian và độ ổn định trên khai cuộc mới cần thí nghiệm bổ sung.

\section*{2. Kế hoạch và tiến độ 8 tuần}

\section*{3. Cấu trúc báo cáo}
Chương 1 giới thiệu bài toán và câu hỏi nghiên cứu. Chương 2 trình bày cơ sở thuật toán. Chương 3 mô tả cách mã hóa bàn cờ, cài đặt tác tử và huấn luyện. Chương 4 trình bày kết quả cùng điều kiện từng phép đo. Chương 5 tổng hợp kết luận và hướng phát triển. Phụ lục~\ref{app:tai-lap} chỉ rõ nguồn số liệu và cách tái lập; các ván minh họa được phân tích ở Chương~\ref{ch:thuc-nghiem}.
